% !TEX program = xelatex
% 공통수학2 · Ⅰ. 도형의 방정식 · 2. 원의 방정식 · 01 원의 방정식 · 2/2차시
\documentclass[11pt,a4paper]{article}
\usepackage{kotex}
\usepackage{iftex}
\ifPDFTeX\else
  \setmainhangulfont{Noto Serif CJK KR}
  \setsanshangulfont{Noto Sans CJK KR}
\fi
\usepackage[margin=2.1cm]{geometry}
\usepackage{parskip}
\usepackage{enumitem}
\usepackage{array}
\usepackage{tabularx}
\usepackage{booktabs}
\usepackage{xcolor}
\usepackage{amssymb}
\usepackage{tikz}
\usepackage[most]{tcolorbox}
\usepackage{fancyhdr}
\usepackage{titlesec}

\definecolor{goldc}{HTML}{B07C22}
\definecolor{goldstrong}{HTML}{8A5F16}
\definecolor{indigoc}{HTML}{2B3B57}
\definecolor{indigosoft}{HTML}{6E82A6}
\definecolor{linec}{HTML}{D6DACB}
\definecolor{surface2}{HTML}{E4E8DC}
\definecolor{notebg}{HTML}{FBF3E3}
\definecolor{inksoft}{HTML}{52604F}

\titleformat{\section}{\normalfont\sffamily\bfseries\large\color{indigoc}}{\thesection}{0.6em}{}
\titlespacing{\section}{0pt}{16pt}{8pt}
\newtcolorbox{ideabox}{enhanced, breakable, colback=white, colframe=goldc, boxrule=0.5pt, leftrule=3pt, arc=2pt, left=10pt, right=10pt, top=8pt, bottom=8pt}
\newtcolorbox{calloutbox}{enhanced, breakable, colback=white, colframe=indigosoft, boxrule=0.5pt, leftrule=3pt, arc=2pt, left=10pt, right=10pt, top=7pt, bottom=7pt}
\newtcolorbox{notebox}{enhanced, breakable, colback=notebg, colframe=goldc, boxrule=0.5pt, leftrule=3pt, arc=2pt, left=10pt, right=10pt, top=7pt, bottom=7pt}
\newtcolorbox{answerbox}{enhanced, breakable, colback=white, colframe=linec, boxrule=0.5pt, arc=2pt, left=10pt, right=10pt, top=7pt, bottom=7pt}
\newcommand{\lessonbanner}[3]{%
  \begin{tcolorbox}[enhanced, colback=indigoc, colframe=indigoc, boxrule=0pt, arc=0pt,
    left=14pt, right=14pt, top=14pt, bottom=10pt, borderline south={2.4pt}{0pt}{goldc}, fontupper=\color{white}]
    {\sffamily\footnotesize\color{white!85!black} #1}\\[4pt]{\sffamily\bfseries\Large #2}\\[3pt]{\sffamily\small\color{white!88!black} #3}
  \end{tcolorbox}}
\pagestyle{fancy}\fancyhf{}\renewcommand{\headrulewidth}{0pt}
\fancyfoot[L]{\footnotesize\color{inksoft} 공통수학2 · 2026학년도 2학기 · 지평선고등학교}
\fancyfoot[R]{\footnotesize\color{inksoft} 교사용 지도안 -- 배포 금지}

\begin{document}
\lessonbanner{공통수학2 · Ⅰ. 도형의 방정식 · 2. 원의 방정식}{01 원의 방정식 --- 2/2차시}{일반형과 세 점을 지나는 원 · 교사용 지도안 (2026학년도 2학기)}
\vspace{10pt}

\noindent\begin{tabularx}{\linewidth}{@{}>{\bfseries\color{indigoc}}p{2.6cm}X@{}}
\toprule
대상 & 고1 · 2개 반 · 반당 13명 \\
차시 & 50분 (도입 10 · 전개 30 · 정리 10) \\
성취기준 & 이차방정식 $x^2+y^2+Ax+By+C=0$이 나타내는 도형을 이해하고, 원의 방정식을 구할 수 있다. \\
선행 학습 & 9차시 --- 원의 방정식(표준형) \\
\bottomrule
\end{tabularx}

\section{학습 목표 \& Big Idea}
\begin{ideabox}
\textbf{\color{goldstrong}영속적 이해 (Big Idea)}\\
같은 원이라도 `표준형'과 `일반형'이라는 두 가지 얼굴을 가진다. 형태는 다르지만 완전제곱식으로 변형하면 서로 오갈 수 있다.
\vspace{4pt}
\small\color{inksoft} 핵심개념: \textbf{동치 변형} \quad\textbullet\quad 핵심개념: \textbf{외접원} \quad\textbullet\quad 일반화: \textbf{같은 대상도 여러 형태의 식으로 표현될 수 있다}
\end{ideabox}
\begin{calloutbox}
\textbf{차시 학습 목표}
\begin{itemize}[leftmargin=1.4em, itemsep=2pt, topsep=4pt]
  \item 세 점을 지나는 원(삼각형의 외접원)의 방정식을 구할 수 있다.
  \item 이차방정식 $x^2+y^2+Ax+By+C=0$이 원을 나타낼 조건을 알고, 표준형으로 변형할 수 있다.
\end{itemize}
\end{calloutbox}

\section{질문의 연결}
\begin{tabularx}{\linewidth}{@{}>{\bfseries\color{indigoc}\small}p{2.4cm}X@{}}
사실적 질문 & 세 점을 지나는 원의 중심은 어떻게 찾을 수 있을까? \\[6pt]
개념적 질문 & 모든 $x^2+y^2+Ax+By+C=0$ 꼴의 식이 원을 나타낼까? 아니라면 언제 아닐까? \\[6pt]
논쟁적 질문 & 같은 대상을 여러 다른 형태로 표현할 수 있다는 것은 수학에서, 그리고 우리 삶에서 어떤 의미가 있을까? \\
\end{tabularx}

\section{수업 흐름}
\begin{tabularx}{\linewidth}{@{}>{\centering\arraybackslash}p{1.6cm}X>{\raggedright\arraybackslash\small}p{3.1cm}@{}}
\toprule
\textbf{단계} & \textbf{활동 \& 발문} & \textbf{ATL / VTR} \\
\midrule
\textcolor{goldstrong}{\textbf{도입}}\newline\footnotesize 10분 &
\textbf{마음공부 (2분)} --- 유무념 대조.\newline
\textbf{생각 열기 (8분)} --- ``한 원 위의 세 점을 알면 원 전체를 알 수 있을까?'' 세 점 A$(1,5)$, B$(2,4)$, C$(-1,1)$을 좌표평면에 찍고 중심의 위치를 어림잡아 보게 한다. &
ATL: 사고(추론)\newline VTR: See-Think-Wonder \\
\addlinespace
\textcolor{goldstrong}{\textbf{전개}}\newline\footnotesize 30분 &
\textbf{개념 형성 (15분)} --- 원의 중심 P$(a,b)$에서 세 점까지의 거리가 모두 같음($\overline{\text{PA}}=\overline{\text{PB}}=\overline{\text{PC}}$)을 이용해 연립방정식을 세우고 예제 2를 함께 풀이 $\to$ 문제 3 짝 활동.\newline
\textbf{적용 (15분)} --- 원의 방정식을 전개하면 $x^2+y^2+Ax+By+C=0$ 꼴이 됨을 확인 $\to$ 완전제곱식으로 변형하여 $A^2+B^2-4C>0$일 때만 원을 나타냄을 유도 $\to$ 문제 4, 5. &
ATT: 촉진자 역할\newline ATL: 사고·변형 \\
\addlinespace
\textcolor{goldstrong}{\textbf{정리}}\newline\footnotesize 10분 &
\textbf{개념 연결 질문 (5분)} --- ``표준형과 일반형, 언제 어느 쪽이 더 편리할까?''\newline
\textbf{성찰 (5분)} --- 감사 일기 1문장 + `원의 방정식' 소단원 마무리 + 다음 차시 예고(원과 직선의 위치 관계). &
마음공부 · 감사 일기 \\
\bottomrule
\end{tabularx}

\begin{calloutbox}
\textbf{지도상의 유의점} --- 세 점을 지나는 원을 구할 때 $\overline{\text{PA}}^2=\overline{\text{PB}}^2$처럼 제곱한 식으로 연립하는 이유(무리식을 피하기 위함)를 짚어 준다. 일반형에서 $A^2+B^2-4C=0$이면 한 점, $<0$이면 도형이 존재하지 않음도 간단히 언급한다.
\end{calloutbox}

\section{형성평가 \& 답안}
\begin{answerbox}
\textbf{문제 3} \quad 세 점 A$(-2,6)$, B$(-5,-3)$, C$(2,-2)$를 지나는 원\\
\textbf{\color{goldstrong}답} \; $(x+2)^2+(y-1)^2=25$
\end{answerbox}
\begin{minipage}[t]{0.48\linewidth}
\begin{answerbox}
\textbf{문제 4} \quad 방정식이 나타내는 도형\\
\small ⑴ $x^2+y^2+4x=0$ \quad ⑵ $x^2+y^2-10x-8y-8=0$\\[4pt]
\textbf{\color{goldstrong}답} \; ⑴ 중심$(-2,0)$, 반지름 $2$ \quad ⑵ 중심$(5,4)$, 반지름 $7$
\end{answerbox}
\end{minipage}\hfill
\begin{minipage}[t]{0.48\linewidth}
\begin{answerbox}
\textbf{문제 5} \quad $x^2+y^2-2x+6y+k=0$이 원이 되도록 하는 $k$의 범위\\[4pt]
\textbf{\color{goldstrong}답} \; $k<10$
\end{answerbox}
\end{minipage}

\section{성취수준 (이 차시 범위)}
\begin{tabularx}{\linewidth}{@{}>{\bfseries\color{goldstrong}}p{0.9cm}X@{}}
\toprule
A & 원의 방정식을 구하고, 일반형 원의 방정식의 그래프를 그릴 수 있다. \\
B & 원의 방정식을 구하고, 일반형을 표준형으로 변형할 수 있다. \\
C & 중심과 반지름이 주어졌을 때 원의 방정식을 구하고 그래프를 그릴 수 있다. \\
D & 표준형 원의 방정식의 그래프를 그릴 수 있다. \\
E & 안내된 절차에 따라 그래프를 그릴 수 있다. \\
\bottomrule
\end{tabularx}

\section{다음 차시}
\begin{calloutbox}
\textbf{02 원과 직선의 위치 관계 --- 1/2차시.} 판별식을 이용하여 원과 직선이 몇 개의 점에서 만나는지 판단하는 방법을 다룬다.
\end{calloutbox}
\end{document}
